\banner{Design project 9: a mulching and tillage plan for a tomato field}\label{sec:proj9}
\textbf{Brief.} A drip-irrigated tomato field on a loam (reference $ET_0=5$ mm\,d$^{-1}$, $a=3.5$ mm\,d$^{-1/2}$) is to be planted after four choices of surface cover. The designer compares the evaporation in the first thirty days after a wetting and the water saved. The evaporation after a wetting is $E=U+a(\sqrt{t}-\sqrt{t_1})$ with $t_1=U/ET_0$; with a cover the value is reduced by $(1-f_c)$.

\begin{table}[!htbp]\centering\small\caption{Thirty-day evaporation from a wetted loam for four surface covers ($ET_0=5$ mm/d, $a=3.5$ mm\,d$^{-1/2}$).}
\begin{tabular}{lrrrrr}\toprule
\hd{Cover} & \hd{Reduction factor} & \hd{$U$ (mm)} & \hd{$E_{30}$ (mm)} & \hd{Saving (mm)} & \hd{Saving (\%)}\\\midrule
Bare soil & 0.00 & 9 & 23.5 & 0.0 & 0\\
Straw 3 t/ha & 0.40 & 8 & 13.6 & 9.8 & 42\\
Straw 6 t/ha & 0.60 & 7 & 8.8 & 14.7 & 62\\
Plastic film & 0.90 & 5 & 2.1 & 21.4 & 91\\\bottomrule\end{tabular}\end{table}

\textbf{Reading the table.} Straw at 6 t\,ha$^{-1}$ saves 15 mm in thirty days, the plastic film 21 mm; over a 110-day season at the same rate, plastic could save as much as 78 mm, equal to the irrigation of three to four rounds. The choice depends on cost, on the rain (film sheds it) and on the disposal of the film.
\fig{g9_evap}{0.7}{Cumulative evaporation after a wetting for four surface covers ($ET_0=5$ mm\,d$^{-1}$, $a=3.5$ mm\,d$^{-1/2}$).}
